\chapter{Answers to all control questions}
\question{1}{What is the difference between a physical core and a logical processor? Where does it appear in our charts?}
A physical core is a hardware execution core. A logical processor is an execution context exposed by the processor and OS to the scheduler. On a processor with simultaneous multithreading, several logical processors may share one physical core's execution resources. On our Apple M3, macOS exposes eight logical processors for eight physical cores: four performance and four efficiency cores. We therefore do not have an additional hardware doubling beyond eight.

In Task 1a, the substantial early gains and much smaller, variable gains beyond 8--16 software threads are consistent with this finite CPU capacity. In Task 1c, the reader and queue limit scaling long before all hardware contexts become useful. The charts show why core count is an upper-resource constraint rather than a guarantee that every algorithm reaches its optimum at exactly that count.

\question{2}{How do we ensure that no line is counted twice or skipped between two file parts?}
We define ownership by the line's starting byte. A worker begins at its nominal boundary only if that boundary follows a newline. Otherwise it advances to the next complete line. At the other end, it finishes any line whose first byte it already owns, even when that line crosses the nominal boundary. The next worker skips that same partial line. These two rules fit together: the earlier worker owns the crossing line, and the later worker owns the next line. We also check all three sign counts against the common reference and verify that the ordered part files reproduce the main file.

\question{3}{How do unsynchronised shared counters lose updates, step by step?}
First, one worker reads a counter value. Before it writes back, another worker reads the same old value. Each computes the old value plus one. Both then write the same new value, so one increment overwrites the effect of the other. Table~\ref{tab:race} shows a concrete interleaving from 100 to 101 instead of 102. This remains possible even when individual reads and writes of the counter fit in a machine word, because the whole increment contains multiple operations. The differing incorrect totals in our three repetitions demonstrate the practical consequence.

\question{4}{Why are per-thread counters faster than shared counters protected by a lock?}
Each worker can increment its own variables without acquiring a shared lock on every number. The hot loops run independently, and the final reduction costs only a few additions per worker after joining. A shared lock serialises the protected block and can force waiting, ownership transfers, and scheduling work. Even an extremely short critical section becomes expensive when we enter it 50 million times. Our eight-thread local-counter measurement is 131.284 ms, versus 1738.592 ms for locking, while both produce the same correct totals.

\question{5}{What does I/O-bound mean for a USB drive and for a network server? Why can their response to concurrency differ?}
I/O-bound means that waiting for input/output, rather than available CPU computation alone, limits progress. For our cold USB run, the process spends much less CPU time than real elapsed time and all reads share the same device. Several threads cannot create additional flash bandwidth. For a server, waiting can concern many independent requests, remote computers, or network round trips, so concurrency can use otherwise idle CPU time while those waits overlap.

The statement that concurrency helps only the network case is a useful simplification, not a universal rule. USB concurrency can also help when it overlaps parsing with reading, hides device latency, or improves the controller's utilisation; our actual speed-up must be read from the measured curve. It eventually encounters a shared device limit. A network service can likewise stop benefiting once its CPU, connection limits, or network bandwidth are saturated.

\question{6}{How does static partitioning differ from taking work from a queue? When is each preferable?}
Static partitioning assigns each worker a fixed range before processing starts. Tasks 1a and 2 use this approach. It has little coordination overhead and is effective when work units are sufficiently uniform. Its weakness is imbalance: a worker that finishes early cannot automatically take a portion of a slower worker's assignment. Equal bytes also do not necessarily mean exactly equal parsing time.

Queue-based scheduling assigns the next available unit to whichever worker finishes and requests more work. Task 1c can therefore balance consumers dynamically when batch costs differ. It also adds a producer, queue operations, synchronisation, and a need for termination/backpressure rules. We prefer static partitioning for cheap, regular work where independence is easy to arrange, and a bounded queue for streaming pipelines or work with variable processing costs. Our measurements show that dynamic scheduling is not automatically faster when the producer itself is the bottleneck.
